% Generated draft: requires mathematical and source review.
\documentclass[11pt]{article}
\usepackage[margin=25mm]{geometry}
\usepackage{fontspec}
\setmainfont{Latin Modern Roman}
\usepackage{amsmath,amssymb,mathtools,unicode-math}
\setmathfont{Latin Modern Math}
\usepackage{enumitem,needspace,xurl,hyperref}
\setlength{\parindent}{0pt}
\setlength{\parskip}{5pt}
\setlength{\emergencystretch}{3em}
\newcommand{\N}{\mathbb{N}}
\newcommand{\Z}{\mathbb{Z}}
\newcommand{\Q}{\mathbb{Q}}
\newcommand{\R}{\mathbb{R}}
\newcommand{\C}{\mathbb{C}}
\newcommand{\F}{\mathbb{F}}
\newcommand{\A}{\mathbb{A}}
\newcommand{\PP}{\mathbb{P}}
\newcommand{\E}{\mathbb{E}}
\newcommand{\eps}{\varepsilon}
\newcommand{\dd}{\,\mathrm{d}}
\begin{document}
\section*{UnsolvedMath ID: 800653}
\textbf{Source problem code:} catalogue entry
\subsection{Definitions and mathematical statement}
The supplied excerpt does not specify the class of spaces under consideration or define its cover conventions. The following makes the usual open-cover interpretation explicit; its applicability to the intended class requires confirmation.

Let \(X\) be a topological space and let \(\N=\{1,2,\ldots\}\). An \(\omega\)-cover of \(X\) is a family \(\mathcal U\) of open subsets of \(X\) such that \(X\notin\mathcal U\) and every finite subset of \(X\) is contained in some member of \(\mathcal U\). A \(\gamma\)-cover is an infinite family of open subsets of \(X\), none equal to \(X\), such that each point of \(X\) belongs to all but finitely many members. Under these conventions, the \(\gamma\)-property means that every open \(\omega\)-cover contains a \(\gamma\)-subcover.

For a sequence \((A_n)_{n\in\N}\) of subsets of \(X\), define
\[
\liminf_{n\to\infty} A_n
=\bigcup_{m\in\N}\bigcap_{n\ge m}A_n.
\]
For \(\mathcal F\subseteq\mathcal P(X)\), let \(L(\mathcal F)\) be the smallest family containing \(\mathcal F\) and closed under this operation. Thus closure includes repeated applications, rather than only one application to sequences from \(\mathcal F\). The space \(X\) has the \(\delta\)-property, or is a \(\delta\)-set, if
\[
\text{for every open }\omega\text{-cover }\mathcal U\text{ of }X,
\qquad X\in L(\mathcal U).
\]

\textbf{Problem.} In the class of spaces intended by the source, is the \(\delta\)-property equivalent to the \(\gamma\)-property? Equivalently, does every space in that class with the \(\delta\)-property have the \(\gamma\)-property?

The source records the implication from the \(\gamma\)-property to the \(\delta\)-property, observing that for a countably infinite \(\gamma\)-cover enumerated without repetitions as \(\{U_n:n\in\N\}\), one has \(X=\liminf_{n\to\infty}U_n\). It presents the converse as open. The intended ambient class, and whether the source restricts \(\omega\)-covers to countable families, remain unresolved by the supplied excerpt; the question above is therefore qualified by those missing specifications.
\subsection{Short English statement}
The question asks whether every \(\delta\)-set must have the \(\gamma\)-property: if repeated sequential lower limits of each open \(\omega\)-cover can produce the whole space, must that cover contain a \(\gamma\)-subcover? The source states the reverse implication and labels the converse open. The intended class of spaces and cover conventions require confirmation.
\subsection{Sources}
\begin{itemize}
\item Local UnsolvedMath registry, permanent ID \(800653\), exact source problem code \texttt{catalogue entry}. The supplied record labels this problem \texttt{open}; this is supplied source status, not an independently verified current assessment.
\item Tsaban, \emph{Selection Principles in Mathematics: A Milestone of Open Problems} (2003), Problem 26. Supplied source URL: \url{https://arxiv.org/abs/math/0312182}. The record reports access on 2026-10-16 and extraction by \texttt{arxiv\_latex}; these are supplied provenance statements, not a claim of access or verification here. The registry lists Tsaban as proposer and 2003 as the proposed year; the excerpt describes the problem as already classical.
\item The excerpt attributes the definition of the \(\delta\)-property to the paper labeled [GN], immediately following its definition of the \(\gamma\)-property, and gives the problem attribution/context as [GN, p.\ 160]. No bibliographic identification of [GN] is included in the supplied record.
\end{itemize}
\end{document}
